\documentclass{beamer}
\mode<presentation> {
% \usetheme{Singapore}
\usetheme{CambridgeUS}
\usecolortheme{seahorse}
\setbeamertemplate{navigation symbols}{}
\setbeamertemplate{footline}[page number]
}
\input{notations}
\setbeamerfont{normal text}{size=\normalsize}
\definecolor{mygreen}{RGB}{0, 128, 0}
\setbeamertemplate{caption}[numbered]
% \usepackage[backend=biber, style=alphabetic]{biblatex}
% \addbibresource{_references.bib}

\usepackage[numbers]{natbib}
\bibliographystyle{plainnat}

% \usepackage{enumitem}
\usepackage{tikz}
\usepackage{wrapfig}

\setbeamertemplate{enumerate item}{\alph{enumi})} % First-level: a) b) c)
\setbeamertemplate{enumerate subitem}{\roman{enumii}.} % Second-level: i. ii. iii.

\newcommand{\<}{\langle}
\renewcommand{\>}{\rangle}
\newcommand{\hy}{\hat{y}}
\newcommand{\bx}{\boldsymbol{x}}
\newcommand{\bw}{\boldsymbol{w}}
\newcommand{\bSigma}{\boldsymbol{\Sigma}}
\newcommand{\btheta}{\boldsymbol{\theta}}
\usepackage{booktabs}
% Set hyperlink colors
\hypersetup{
    colorlinks=true,
    linkcolor=mygreen,   % Color of internal links (sections, pages, etc.)
    citecolor=mygreen,   % Color of citation links (bibliography)
    urlcolor=mygreen     % Color of external links (URLs)
}
\setbeamercolor{footnote mark}{fg=mygreen}
% \renewcommand*{\thefootnote}{\textcolor{orange}{\arabic{footnote}}}


\AtBeginDocument{\usebeamerfont{normal text}}


\begin{document}
\begin{frame}
    \title{Underdamped Langevin Algorithm and It's Applications}
    \author{Dmitrii Polozkov\\
    (under supervision of Dr. J-F. Jabir)}
    \begin{center}
        \includegraphics[width=0.5\textwidth]{figs/defend ba/hse logo.png} % Replace with your figure file name
    \end{center}
    \date{30.03.2026}
    \maketitle
\end{frame} 

\begin{frame}{Motivation}

    \begin{center}
        Many modern optimizers can be interpreted as \textbf{discretizations of continuous dynamics}.
    \end{center}

    \bigskip

    \begin{itemize}
        \item Gradient descent = discretization of \textbf{gradient flow}
        \item SGD = discretization of \textbf{overdamped Langevin dynamics}
    \end{itemize}

    \bigskip

    \begin{block}{Design principle}
        Construct new optimization algorithms by choosing a continuous dynamic and discretizing it appropriately
    \end{block}

    \bigskip

    \begin{center}
        Let's work with \textbf{underdampled Langevin dynamics}
    \end{center}
    
\end{frame}

\begin{frame}{SGD as discretization of overdamped Langevin}

    \begin{itemize}
        \item Start from the continuous dynamic (overdamped Langevin):
        \[
        d\theta_t=-\nabla V(\theta_t)dt+\sigma dW_t
        \]
        
        \bigskip

        \item Apply Euler--Maruyama discretization with step $\varepsilon_n$:
        \[
        \theta_{n+1}
        =\theta_n-\varepsilon_n\nabla V(\theta_n)
        +\sqrt{\varepsilon_n}\sigma\xi_n,
        \qquad \xi_n\sim\mathcal N(0,I_d)
        \]

        \bigskip

        \item This matches SGD:
        \[
        \theta_{n+1}=\theta_n-\varepsilon_n\nabla \hat V(\theta_n)
        \]
        
        \item The estimator $\nabla \hat V(\theta_n)$ satisfies
        \[
        \mathbb E[\nabla \hat V(\theta_n)\mid \theta_n]=\nabla V(\theta_n)
        \]
        but has non-zero variance, introducing \textbf{stochastic noise} in each update
    \end{itemize}
\end{frame}

\begin{frame}{Underdamped dynamics and the small-$\gamma$ regime}
    \begin{block}{Underdamped Langevin dynamics}
        \[
        \begin{cases}
            dX_t^\gamma = V_t^\gamma\,dt,\\[0.3em]
            dV_t^\gamma = -\gamma^{-1}\bigl(\nabla f(X_t^\gamma)+V_t^\gamma\bigr)\,dt
            +\gamma^{-1}\sigma\,dW_t
        \end{cases}
        \qquad t\ge 0
        \]
    \end{block}
    \begin{itemize}
        \item The velocity variable reacts on a much faster time scale than the position because of $\gamma$
        \item When $\gamma$ is small, inertia is quickly forgotten, so the motion is driven mainly by the current landscape and noise\\
        
        \medskip 
        $\Rightarrow$ Position $X_t^\gamma$ should then behave like the overdamped process $X_t^0$
    \end{itemize}
    \begin{block}{Overdamped limit}
        \[
            dX_t^0=-\nabla f(X_t^0)\,dt+\sigma\,dW_t
        \]
    \end{block}
\end{frame}

\begin{frame}{Continuous convergence results proved so far}
    Assume $\nabla f$ is globally Lipschitz and the initial condition has finite moments.
    \bigskip
    \begin{block}{\textbf{Theorem} \textcolor{mygreen}{Time-marginal mean-square convergence}}
        For every finite $T$ and every integer $p\ge 1$,
        \[
        \max_{t\in[0,T]}\E\bigl[|X_t^\gamma-X_t^0|^{2p}\bigr]=\bigO(\gamma^p).
        \]
    \end{block}
    Let's study the \textcolor{mygreen}{time-marginal} bound for the discrete schemes
    \bigskip
    \begin{block}{\textbf{Proposition} Max-in-time estimate}
        Also,
        \[
        \E\Big[\max_{t\in[0,T]}|X_t^\gamma-X_t^0|^2\Big]=\bigO(\gamma^{1/2}).
        \]
        (\textit{still under clarification})
    \end{block}
\end{frame}


\begin{frame}
\centering
\begin{tikzpicture}[>=stealth]

\node[draw] (root) at (0,0) {\Large Langevin equation \normalsize};

\node[draw, align=center] (OLD) at (-3, -2) {
    overdamped Langevin\\
    dynamics\\
    \[
    \small
    dX_t^0=-\nabla f(X_t^0)\,dt+\sigma\,dW_t
    \normalsize
    \]
};

\node[draw, align=center] (ULD) at (3, -2) {
    underdamped Langevin\\
    dynamics\\
    \[
    \tiny
    \begin{cases}
        dX_t^\gamma=V_t^\gamma\,dt\\
        dV_t^\gamma=
        -\gamma^{-1}\bigl(\nabla f(X_t^\gamma)+V_t^\gamma\bigr)\,dt
        +\gamma^{-1}\sigma\,dW_t
    \end{cases}
    \normalsize
    \]
};

% No arrows from root

% Arrow: underdamped -> overdamped with gamma -> 0
\draw[->, thick, draw=mygreen]
  (ULD) -- (OLD)
  node[midway, above, text=mygreen] {\tiny$\gamma \to 0$};

\node[draw, align=center] (OLA) at (-3, -5) {
    SGD\\
    aka overdamped Langevin algorithm\\
    \[
    \small
    \theta_{n+1} = \theta_n-\varepsilon_n\nabla\hat{V}(\theta_n)
    \normalsize
    \]
};

\node[draw=cyan, text=cyan, align=center] (ULA) at ( 3, -5) {
    what numerical scheme\\
    to use for\\
    discretization of\\
    underdamped?\\[0.2em]
    maybe SGD too?
};

\draw[->, dashed, thick]
  (OLD) -- (OLA)
  node[midway, left] {\small Euler-Maruyama\normalsize};

\draw[->, dashed, thick, draw=cyan]
  (ULD) -- (ULA);

\end{tikzpicture}
\end{frame}

\begin{frame}{Why plain SGD / Euler--Maruyama is not enough}
    Applying Euler--Maruyama to the underdamped system gives
    \[
    \begin{cases}
        X_{n+1}=X_n+hV_n,\\
        V_{n+1}=V_n-\dfrac{h}{\gamma}\bigl(\nabla f(X_n)+V_n\bigr)
        +\dfrac{\sigma\sqrt h}{\gamma}\xi_n
    \end{cases}
    \]
    \begin{itemize}
        \item Velocity has \textbf{exponential relaxation} at rate $1/\gamma$
        \item As $\gamma\downarrow 0$, this becomes \textbf{stiff} $\Rightarrow$ \textbf{instability}
    \end{itemize}
    \bigskip
    \begin{block}{Conclusion}
        The discrete scheme must capture the fast Ornstein--Uhlenbeck relaxation of the velocity equation.
    \end{block}
\end{frame}

% \begin{frame}{Design of the scheme: exponential integration of the OU part}
%     Freeze the force on one step: $\nabla f(X_t^\gamma)\approx \nabla f(X_n)$ on $[t_n,t_{n+1}]$, and solve the linear velocity equation exactly.
%     \bigskip
%     Let $a=e^{-h/\gamma}$. Then the proposed one-step update is
%     \[
%     V_{n+1}=aV_n-(1-a)\nabla f(X_n)+\eta_n^{(v)},
%     \]
%     \[
%     X_{n+1}=X_n+\gamma(1-a)V_n-\bigl(h-\gamma(1-a)\bigr)\nabla f(X_n)+\eta_n^{(x)},
%     \]
%     where $(\eta_n^{(x)},\eta_n^{(v)})$ is the exact Gaussian noise generated by the frozen linear system.
%     \bigskip
%     \begin{itemize}
%         \item The linear drift in the velocity is solved exactly over each step
%         \item The exponential form captures the rapid decay of velocity (relaxation)
%         \item The only approximation comes from keeping $\nabla f(X_t^\gamma)$ fixed within each timestep
%     \end{itemize}
% \end{frame}

\begin{frame}{Design of the scheme: exact solution of frozen OU dynamics}

    On each interval $t\in[t_n, t_{n+1}]$, freeze the force:
    \[
    \nabla f(X_t^\gamma)\approx \nabla f(X_{t_n}).
    \]

    Then the velocity equation becomes a linear SDE (Ornstein--Uhlenbeck):
    \[
    dV_t = -\tfrac{1}{\gamma}V_t\,dt - \tfrac{1}{\gamma}\nabla f(X_n)\,dt + \tfrac{\sigma}{\gamma}dW_t.
    \]

    This equation admits an explicit solution:
    \[
    V_{n+1} = e^{-h/\gamma}V_n - (1-e^{-h/\gamma})\nabla f(X_n)
    + \int_{t_n}^{t_{n+1}} e^{-(t_{n+1}-s)/\gamma}\tfrac{\sigma}{\gamma}\,dW_s.
    \]

    Denote $a=e^{-h/\gamma}$. Then
    \[
    \boxed{
    V_{n+1}=aV_n-(1-a)\nabla f(X_n)+\eta_n^{(v)}.
    }
    \]

    The position is obtained by integrating $dX_t = V_t dt$:
    \[
    \boxed{
    X_{n+1}=X_n+\gamma(1-a)V_n-\bigl(h-\gamma(1-a)\bigr)\nabla f(X_n)+\eta_n^{(x)}.
    }
    \]

\end{frame}

\begin{frame}{Gaussian noise: exact law}

    The noise terms are stochastic convolutions of Brownian motion:
    \[
    \eta_n^{(v)}=\int_{t_n}^{t_{n+1}} e^{-(t_{n+1}-s)/\gamma}\tfrac{\sigma}{\gamma}\,dW_s,
    \]
    \[
    \eta_n^{(x)}=\int_{t_n}^{t_{n+1}}\int_{t_n}^{t} e^{-(t-s)/\gamma}\tfrac{\sigma}{\gamma}\,dW_s\,dt.
    \]

    They are jointly Gaussian with zero mean and explicit covariance:
    \[
    \mathbb E[\eta_n^{(v)}(\eta_n^{(v)})^\top]
    = \tfrac{\sigma^2}{2\gamma}(1-a^2)I,
    \]
    \[
    \mathbb E[\eta_n^{(x)}(\eta_n^{(x)})^\top]
    = \sigma^2\Big(h - 2\gamma(1-a) + \tfrac{\gamma}{2}(1-a^2)\Big)I,
    \]
    \[
    \mathbb E[\eta_n^{(x)}(\eta_n^{(v)})^\top]
    = \tfrac{\sigma^2}{2}(1-a)^2 I.
    \]

    Hence $(\eta_n^{(x)},\eta_n^{(v)})$ is sampled from a known Gaussian law.

\end{frame}

\begin{frame}{Key properties of the scheme}

    \begin{itemize}
        \item \textbf{Exact treatment of stiffness:} the linear relaxation in $V_t$ is solved analytically, avoiding instability of Euler--Maruyama for small $\gamma$
        
        \item \textbf{Correct scaling in $\gamma$:} all coefficients and noise moments remain uniformly bounded as $\gamma\to 0$
        
        \item \textbf{Minimal approximation:} the only discretization error comes from freezing $\nabla f$, which is assumed regular enough
    \end{itemize}

\end{frame}

\begin{frame}{Discrete convergence setup}
    Goal: compare the discrete underdamped position $\bar X_t^\gamma$ and the discrete overdamped position $\bar X_t^0$.
    \bigskip
    \begin{block}{Triangle inequality decomposition}
        \[
        \sup_{t\in[0,T]}\E|\bar X_t^\gamma-\bar X_t^0|^2
        \le
        C\sup_{t\in[0,T]}\E|\bar X_t^\gamma-X_t^\gamma|^2
        \]
        \[
        \qquad\qquad
        +C\sup_{t\in[0,T]}\E|X_t^\gamma-X_t^0|^2
        +C\sup_{t\in[0,T]}\E|X_t^0-\bar X_t^0|^2
        \]
    \end{block}
    \begin{itemize}
        \item Middle term is already known:
        \[
        \sup_{t\in[0,T]}\E|X_t^\gamma-X_t^0|^2=\bigO(\gamma)
        \]
        \item Last term is the standard overdamped discretization error
        \item Only tricky term:
        \[
        \sup_{t\in[0,T]}\E|\bar X_t^\gamma-X_t^\gamma|^2
        \]
    \end{itemize}
\end{frame}

\begin{frame}{Current analytical target and what is left open}
    Desired bound for the exponential scheme:
    \[
    \sup_{t\in[0,T]}\E|\bar X_t^\gamma-X_t^\gamma|^2\le Ch^\alpha
    \]
    with a constant $C$ \textbf{independent of $\gamma$}.
    \bigskip
    \begin{itemize}
        \item If $\alpha=1$, then under $h\lesssim\gamma$ this would preserve the continuous $O(\gamma)$ rate at the discrete level
        \item This is the main unresolved step in the discrete theory
    \end{itemize}

    \bigskip
    
    \textit{In future, the discrete analogue of max-in-time continuous rate is also established in a similar form}
\end{frame}

\begin{frame}{Numerical test: quadratic benchmark in dimension $d=1$}
    Use the simple potential
    \[
    f(x)=\tfrac{\lambda}{2}x^2,\qquad \nabla f(x)=\lambda x
    \]
    so the continuous underdamped dynamics is linear:
    \[
    \begin{cases}
    dX_t^\gamma=V_t^\gamma\,dt,\\
    dV_t^\gamma=-\gamma^{-1}(\lambda X_t^\gamma+V_t^\gamma)\,dt
    +\gamma^{-1}\sigma\,dW_t
    \end{cases}
    \]
    \begin{itemize}
        \item This case admits an exact Gaussian solution
        \item Hence $X_t^\gamma$ can be computed exactly and compared to $\bar X_t^\gamma$
        \item For fixed $\gamma$ and several step sizes $h$, measure
        \[
        \max_{t_n\le T}\E|\bar X_{t_n}^\gamma-X_{t_n}^\gamma|^2
        \]
        \item Then test whether the observed law is
        \[
        \max_{t_n\le T}\E|\bar X_{t_n}^\gamma-X_{t_n}^\gamma|^2
        \approx C(\gamma)h^\alpha
        \]
        with $C(\gamma)$ stable as $\gamma\downarrow0$
    \end{itemize}
\end{frame}

\begin{frame}{Current research status}
    \begin{itemize}
        \item \textcolor{mygreen}{Continuous side:}
        \begin{itemize}
            \item time-marginal convergence
            \[
            \max_{t\in[0,T]}\E|X_t^\gamma-X_t^0|^2=\bigO(\gamma)
            \]
            is established
            \item the max-in-time rate
            \textcolor{lightgray}{
            \[
            \E\Big[\max_{t\in[0,T]}|X_t^\gamma-X_t^0|^2\Big]=\bigO(\gamma^{1/2})
            \]
            }
            is still being checked and refined
        \end{itemize}
        \bigskip
        \item \textcolor{mygreen}{Discrete side:}
        \begin{itemize}
            \item exponential discretization for the underdamped dynamics is proposed
            \item main open goal is the uniform-in-$\gamma$ bound
            \[
            \sup_{t\in[0,T]}\E|\bar X_t^\gamma-X_t^\gamma|^2\le Ch^\alpha
            \]
            \item parallel numerical tests are used to assess whether $\alpha=1$ and whether $C$ remains bounded as $\gamma\downarrow0$
        \end{itemize}
    \end{itemize}
\end{frame}

% \begin{frame}
% \frametitle{References}
    % \printbibliography
% \end{frame}

% \begin{frame}[allowframebreaks]{References}
%     \bibliography{_references}
% \end{frame}

\begin{frame}
    \begin{center}
        \huge Ready to answer your questions
    \end{center}
\end{frame}

\end{document}